\documentclass[a4paper,11pt]{jsarticle}
\usepackage{amsmath,amssymb}

\begin{document}

\title{分数型漸化式で与えられる数列の極限値}
\date{}
\maketitle

\section*{問題}
次の漸化式で定められる数列 $\{a_n\}$ の極限値を求めよ。
\[
a_1 = 1, \quad a_{n+1} = 1 + \frac{1}{1 + a_n} \quad (n = 1, 2, 3, \dots)
\]

\section*{解答}
\subsection*{1. 特性方程式と漸化式の変形}
漸化式 $a_{n+1} = f(a_n)$ における特性方程式 $\alpha = 1 + \frac{1}{1 + \alpha}$ を解く。
両辺に $1 + \alpha$ を掛けると、
\[
\alpha(1 + \alpha) = (1 + \alpha) + 1 \implies \alpha^2 + \alpha = \alpha + 2 \implies \alpha^2 = 2
\]
よって、$\alpha = \pm\sqrt{2}$ である。

ここで、漸化式の両辺から $\sqrt{2}$ を引くと、
\begin{align*}
a_{n+1} - \sqrt{2} &= 1 + \frac{1}{1 + a_n} - \sqrt{2} \\
&= \frac{1 + a_n + 1 - \sqrt{2}(1 + a_n)}{1 + a_n} \\
&= \frac{(1-\sqrt{2})a_n + 2 - \sqrt{2}}{1 + a_n} \\
&= \frac{(1-\sqrt{2})(a_n - \sqrt{2})}{1 + a_n} \quad \cdots (1)
\end{align*}
同様に、両辺に $\sqrt{2}$ を足すと、
\begin{align*}
a_{n+1} + \sqrt{2} &= 1 + \frac{1}{1 + a_n} + \sqrt{2} \\
&= \frac{1 + a_n + 1 + \sqrt{2}(1 + a_n)}{1 + a_n} \\
&= \frac{(1+\sqrt{2})a_n + 2 + \sqrt{2}}{1 + a_n} \\
&= \frac{(1+\sqrt{2})(a_n + \sqrt{2})}{1 + a_n} \quad \cdots (2)
\end{align*}

\subsection*{2. 等比数列への帰着}
すべての自然数 $n$ において $a_n > 0$ は明らかであるため、$a_n + \sqrt{2} \neq 0$ である。(1)を(2)で割ると、
\[
\frac{a_{n+1} - \sqrt{2}}{a_{n+1} + \sqrt{2}} = \frac{1 - \sqrt{2}}{1 + \sqrt{2}} \cdot \frac{a_n - \sqrt{2}}{a_n + \sqrt{2}}
\]
ここで、数列 $\{b_n\}$ を $b_n = \frac{a_n - \sqrt{2}}{a_n + \sqrt{2}}$ とおくと、
\[
b_{n+1} = \frac{1 - \sqrt{2}}{1 + \sqrt{2}} b_n
\]
となり、$\{b_n\}$ は公比 $\frac{1 - \sqrt{2}}{1 + \sqrt{2}}$ の等比数列となる。
初項 $b_1$ は、
\[
b_1 = \frac{a_1 - \sqrt{2}}{a_1 + \sqrt{2}} = \frac{1 - \sqrt{2}}{1 + \sqrt{2}}
\]
であるから、一般項 $b_n$ は、
\[
b_n = b_1 \cdot \left( \frac{1 - \sqrt{2}}{1 + \sqrt{2}} \right)^{n-1} = \left( \frac{1 - \sqrt{2}}{1 + \sqrt{2}} \right)^n
\]
（※ $\frac{1 - \sqrt{2}}{1 + \sqrt{2}}$ は分母を有理化すると $2\sqrt{2} - 3$ になるため、$b_n = (2\sqrt{2} - 3)^n$ と書くこともできる。）

\subsection*{3. 一般項 $a_n$ の導出}
$b_n = \frac{a_n - \sqrt{2}}{a_n + \sqrt{2}}$ を $a_n$ について解く。
\begin{align*}
b_n (a_n + \sqrt{2}) &= a_n - \sqrt{2} \\
a_n(1 - b_n) &= \sqrt{2}(1 + b_n) \\
a_n &= \sqrt{2} \frac{1 + b_n}{1 - b_n}
\end{align*}
これに $b_n = (2\sqrt{2} - 3)^n$ を代入して、一般項を得る。
\[
a_n = \sqrt{2} \frac{1 + (2\sqrt{2} - 3)^n}{1 - (2\sqrt{2} - 3)^n}
\]

\subsection*{4. 極限値の計算}
次に、$n \to \infty$ のときの $a_n$ の極限を求める。
ここで公比 $r = 2\sqrt{2} - 3$ について考えると、$\sqrt{8} \approx 2.828$ より $2\sqrt{2} - 3 \approx -0.172$ であるため、
$|2\sqrt{2} - 3| < 1$ となる。
したがって、
\[
\lim_{n \to \infty} (2\sqrt{2} - 3)^n = 0
\]
である。ゆえに、求める極限値は、
\[
\lim_{n \to \infty} a_n = \lim_{n \to \infty} \sqrt{2} \frac{1 + (2\sqrt{2} - 3)^n}{1 - (2\sqrt{2} - 3)^n} = \sqrt{2} \frac{1 + 0}{1 - 0} = \sqrt{2}
\]

\vspace{1em}
\noindent\textbf{【答え】}\\
一般項: $\displaystyle a_n = \sqrt{2} \frac{1 + (2\sqrt{2} - 3)^n}{1 - (2\sqrt{2} - 3)^n}$ \\[1em]
極限値: $\displaystyle \lim_{n \to \infty} a_n = \sqrt{2}$


% 以下は別解


\newpage

\section*{別解(コーシー列による解法)}

\vspace{1em}
\noindent
$\displaystyle |a_n - a_m| \leqq \frac{|a_{m-1} - a_{n-1}|}{(1 + a_{m-1})(1 + a_{n-1})}$
と計算できるので、$(1 + a_{m-1})(1 + a_{n-1})$ の評価をする.
\vspace{1em}

まず $1 \leqq a_n \leqq \dfrac{3}{2}$ であることを帰納法により示す.

\begin{itemize}
    \item $a_1 = 1$ である.
    \item $1 \leqq a_{n-1} \leqq \dfrac{3}{2}$ とすると $\dfrac{2}{5} \leqq \dfrac{1}{1 + a_{n-1}} \leqq \dfrac{1}{2}$ であるため、特に $0 \leqq \dfrac{1}{1 + a_{n-1}} \leqq \dfrac{1}{2}$ である。よって
    \[
    1 \leqq a_n \leqq \frac{3}{2}
    \]
    となり示された.
\end{itemize}

これより $(1 + a_{m-1})(1 + a_{n-1}) \geqq 4$ がわかる.

正の実数 $\varepsilon$ を任意に取る.

$n > m$ とする.

\begin{align*}
|a_n - a_m| &= \left| \left( 1 + \frac{1}{1 + a_{n-1}} \right) - \left( 1 + \frac{1}{1 + a_{m-1}} \right) \right| \\
&= \left| \frac{a_{m-1} - a_{n-1}}{(1 + a_{m-1})(1 + a_{n-1})} \right| \\
&\leqq \frac{|a_{n-1} - a_{m-1}|}{4} \\
&\leqq \frac{|a_{n-m+1} - a_1|}{4^{m-1}} \\
&\leqq \frac{1}{2 \cdot 4^{m-1}}
\end{align*}

$m \to \infty$ で $\dfrac{1}{2 \cdot 4^{m-1}} \to 0$ であるため、$\{a_n\}$ はコーシー列である。よって、収束するのでその値を $\lim_{n\to\infty} a_n = \alpha$ とおく.

\begin{align*}
\alpha &= \lim_{n\to\infty} a_{n} \\
&= \lim_{n\to\infty} \left( 1 + \frac{1}{1 + a_{n-1}} \right) \\
&= 1 + \frac{1}{1 + \alpha}
\end{align*}

である。整理すると
\[
\alpha^2 - 2 = 0
\]
であるため $\alpha = \pm\sqrt{2}$ である.

$a_n > 0$ であるため $\alpha > 0$ である. よって $\alpha = \sqrt{2}$ である.

\end{document}
